\subsection{Gerenciamento de Infraestruturas de Cloud Computing}

A ascensão da computação em nuvem consolidou o modelo de entrega de recursos como serviço (\textit{Everything as a Service - XaaS}), exigindo infraestruturas capazes de suportar flutuações drásticas de demanda com latência mínima. Nesse contexto, a Computação Autônoma atua como o motor de orquestração essencial para manter a viabilidade econômica e operacional dos provedores de nuvem. O gerenciamento de um \textit{Data Center} moderno envolve milhares de servidores físicos e milhões de máquinas virtuais (VMs) ou containers, onde decisões sobre alocação de recursos precisam ser tomadas em milissegundos \cite{lalanda2013}.

A propriedade de auto-otimização é a mais visível nesta aplicação. Através do ciclo MAPE-K, gerentes autônomos monitoram métricas de desempenho como utilização de CPU, consumo de memória e tráfego de rede. Quando a carga de trabalho de uma aplicação atinge um limiar crítico, o sistema inicia um processo de \textit{auto-scaling}, provisionando novos recursos sem intervenção humana para garantir o cumprimento dos Acordos de Nível de Serviço (SLAs). Inversamente, em períodos de baixa demanda, o sistema consolida as cargas de trabalho em um número menor de servidores físicos e desliga os equipamentos ociosos, resultando em uma economia significativa de energia e redução de custos operacionais \cite{kephart2003}.

Além da eficiência energética, a auto-cura é vital para a resiliência da nuvem. Em sistemas de larga escala, falhas de hardware são eventos estatisticamente certos. Gerentes autônomos detectam a degradação de componentes ou a falha total de um nó e, automaticamente, migram as VMs afetadas para servidores saudáveis, reconfigurando as tabelas de roteamento e balanceadores de carga de forma transparente para o usuário final \cite{correa2007}. Este nível de automação permite que a infraestrutura mantenha a alta disponibilidade mesmo sob condições de estresse ou falhas parciais.

\subsection{Sistemas Pervasivos e Internet das Coisas (IoT)}

Se no modelo de nuvem a complexidade reside na densidade de recursos, nos sistemas pervasivos e na Internet das Coisas (IoT), o desafio reside na heterogeneidade e na volatilidade das conexões. Cidades inteligentes, redes elétricas inteligentes (\textit{Smart Grids}) e ambientes hospitalares conectados dependem de uma miríade de dispositivos sensores e atuadores com capacidades computacionais limitadas e fontes de energia finitas \cite{sterritt2005}.

A auto-configuração é o pilar fundamental para a escalabilidade da IoT. Em uma \textit{Smart City}, novos sensores de monitoramento de tráfego ou qualidade do ar são adicionados à rede diariamente. Um sistema autônomo permite que esses dispositivos, ao serem ligados, descubram vizinhos, estabeleçam protocolos de comunicação seguros e registrem seus serviços no diretório central de forma plug-and-play, mimetizando a maneira como novas células se integram ao corpo humano \cite{kephart2003}. Sem essa capacidade, o custo de instalação e configuração individual de milhões de dispositivos tornaria esses projetos financeiramente proibitivos.

Adicionalmente, a auto-proteção ganha uma nova dimensão em sistemas pervasivos. Devido à natureza distribuída e, muitas vezes, fisicamente exposta dos dispositivos IoT, eles são alvos frequentes de ataques de negação de serviço distribuído (DDoS) ou tentativas de infiltração. Gerentes autônomos de segurança utilizam modelos de inteligência artificial para identificar padrões de tráfego anômalos em tempo real. Ao detectar um comportamento suspeito em um grupo de sensores, o sistema pode isolar logicamente esses componentes do restante da rede (\textit{quarentena}), impedindo a propagação de \textit{malware} ou a exfiltração de dados sensíveis antes mesmo que os administradores humanos tomem conhecimento do incidente \cite{lalanda2013}.